\documentsclass[11pt]{article}
\usepackage{parskip}
\usepackage{color}

\title{Corn chowder}
\author{Nathaniel Zuk}

\begin{document}
\maketitle

\textcolor{green}{(12/10/2017)}

\textcolor{green}{Based on recipe from Professional Chef, p. 434.}

\textbf{Ingredients:}

\begin{enumerate}
\item Onions = 123 g
\item Celery = 182 g
\item Peppers (Green + Red) = 351 g
\item Potatoes (cubed) = 1068 g
\item Chicken stock = 2 L
\item All-purpose flour = 94 g
\item Butter = 2 oz \textcolor{red}{(added 2 Tbsp more butter because this wasn't enough once the flour was added)}
\item Corn kernels = 2 lb
\item Bay leaf = 1
\item Heavy cream = 8 fl oz
\item Milk = 8 fl oz
\item Salt ~ 2 tsp
\item Ground pepper ~ 1 tsp
\item Red hot sauce = 2 tsp
\item Worcestershire sauce = 2 tsp
\end{enumerate}

In a large pot, sautee the onions, celery, and peppers in butter until softened.  Then add flour and cook to make a white roux, ~3 min \textcolor{red}{(although it took me less time, the flour cooked very quickly)}.  Remove the pot from the heat and add 1/2 the stock.  Stir until combined, then return to heat and stir until there are no lumps (from the flour).  Do this a second time with the other 1/2 of the stock.  Simmer for 30 minutes, stirring periodically, until it thickens.  Puree half of the corn and add it to the soup with the potatoes.  Then add the rest of the corn and bay leaf and simmer until the corn and potatoes are tender (20 min).  Add cream and milk to the soup and continue to simmer for 10 min.  Add salt, pepper, red hot and worcestershire sauce after removing from heat.

\textcolor{red}{I think the cream and milk are unnecessary, the soup would be good without it.}

\textbf{Serving size is ~400 g.}

\end{document}